% Generated draft: requires mathematical and source review.
\documentclass[11pt]{article}
\usepackage[margin=25mm]{geometry}
\usepackage{fontspec}
\setmainfont{Latin Modern Roman}
\usepackage{amsmath,amssymb,mathtools,unicode-math}
\setmathfont{Latin Modern Math}
\usepackage{enumitem,needspace,xurl,hyperref}
\setlength{\parindent}{0pt}
\setlength{\parskip}{5pt}
\setlength{\emergencystretch}{3em}
\newcommand{\N}{\mathbb{N}}
\newcommand{\Z}{\mathbb{Z}}
\newcommand{\Q}{\mathbb{Q}}
\newcommand{\R}{\mathbb{R}}
\newcommand{\C}{\mathbb{C}}
\newcommand{\F}{\mathbb{F}}
\newcommand{\A}{\mathbb{A}}
\newcommand{\PP}{\mathbb{P}}
\newcommand{\E}{\mathbb{E}}
\newcommand{\eps}{\varepsilon}
\newcommand{\dd}{\,\mathrm{d}}
\begin{document}
\section*{UnsolvedMath ID: 30001418}
\textbf{Source problem code:} catalogue entry
\subsection{Definitions and mathematical statement}
A Lipschitz interface is a separating surface that can locally be represented, after a rigid change of coordinates, as the graph of a Lipschitz function. Time-harmonic Maxwell equations describe electromagnetic fields with a fixed temporal frequency. A finite element approximation uses finite-dimensional spaces associated with a mesh to approximate the unknown fields.

\textbf{Source-level problem.} For time-harmonic Maxwell equations with material coefficients that change sign across a general Lipschitz interface, establish convergence of the standard finite element schemes. The supplied background specifically identifies sign-changing permittivity as the outstanding setting, whereas the statement refers more broadly to sign-changing material coefficients; the intended scope of the other coefficients is unresolved.

Schematically, if \(U\) denotes the exact solution of the intended Maxwell problem and \(U_h\) its finite element approximation on a mesh of size \(h\), the requested convergence would take a form such as
\[
\lim_{h\to 0}\|U-U_h\|_X=0,
\]
where the solution space \(X\), its norm, the admissible meshes, and the discrete scheme must be specified. This display only explains the meaning of a possible convergence target; the supplied record does not select this topology or provide hypotheses under which such a limit is asserted.

\textbf{Unresolved qualifications.} The record does not specify the spatial domain or dimension, the Maxwell formulation, the frequency regime, the forcing, boundary or radiation conditions, or transmission conditions. It also does not specify whether the material coefficients are scalar or tensor-valued, real or complex, their regularity and bounds, the meaning of sign change in the intended coefficient class, or restrictions on material contrast. The finite element family, interface treatment, mesh assumptions, and convergence notion are likewise unspecified. Solvability, uniqueness, resonance exclusions, and any required stability hypotheses must therefore be resolved before this request can be expressed as a fully quantified convergence theorem. In particular, the source does not support a universal assertion for every sign-changing coefficient or every standard finite element scheme.
\subsection{Short English statement}
Establish convergence of the usual finite element methods for time-harmonic Maxwell equations across a general Lipschitz interface with sign-changing material coefficients, particularly permittivity. The supplied record leaves the precise equations, admissible coefficients, boundary conditions, numerical methods, and convergence topology unresolved. Its partially solved label and literature assessment are supplied dated curation, not independently verified mathematical conclusions.
\subsection{Sources}
\begin{itemize}
\item Supplied local UnsolvedMath registry record, permanent ID \(30001418\), source problem code \texttt{catalogue entry}, title \emph{Finite Elements for Sign-Changing Maxwell Interfaces}, in the collection \emph{Oberwolfach Reports Open Problems}. The record gives the proposed year as 2010 and names no individual proposer.
\item Supplied report provenance: \emph{Computational Electromagnetism and Acoustics} (2010), Oberwolfach Reports, DOI 10.4171/owr/2010/10, pp. 415--519. \url{https://doi.org/10.4171/owr/2010/10}. The supplied background describes Oberwolfach Reports as documenting mathematical workshops and their open-problem discussions.
\item Supplied literature reference: Chesnel--Ciarlet, 2011. \url{https://doi.org/10.1137/100810903}. No article title or precise result is supplied.
\item Supplied literature reference: Chen--Du--Zou, 2000. \url{https://doi.org/10.1137/S0036142998349977}. No article title or precise result is supplied.
\item Supplied literature reference: Oberwolfach Report 10/2010. \url{https://doi.org/10.4171/owr/2010/10}.
\item Literature assessment dated exactly \texttt{2026-08-21}: the supplied curation labels the problem partially solved and reports established finite element convergence in several discontinuous-coefficient and metamaterial settings, while describing the usual finite element theory for a general Lipschitz interface with sign-changing permittivity as unresolved. This is supplied dated curation derived from the OWR material, not independently verified; time-sensitive status and cited claims require independent verification.
\item The registry marks the statement as \texttt{exact} and reports that its mathematical content was checked against the official report, with notation, context, and question form normalized without changing the source-level claim. This verification description is supplied metadata; no independent source inspection is asserted here.
\end{itemize}
\end{document}
